\documentclass[12pt,a4paper]{article}

\usepackage[brazilian]{babel}

\usepackage{amsmath}
\usepackage{amssymb}
\usepackage{enumitem}
\usepackage{etoolbox}
\usepackage{tocloft}
\usepackage{tikz}

\setlength{\cftsecnumwidth}{3.2cm}
\setlength{\cftsubsecnumwidth}{1.5cm}

\renewcommand{\thesection}{Capítulo \arabic{section} -}
\renewcommand{\thesubsection}{\arabic{section}.\arabic{subsection}}
\pretocmd{\subsection}{\clearpage}{}{}

\title{
    \Huge \textbf{Um Curso de Cálculo Volume 1} \\
    \vspace{0.3cm}
    \Large 6\textordfeminine{} Edição (Guidorizzi) \\
    \vspace{0.5cm}
    \large \textit{Resoluções Completas e Passo a Passo do Capítulo 4 Extensões do Conceito de Limite}
}
\author{Rogério Campos}

\begin{document}

\maketitle
\newpage

\tableofcontents
\newpage

\setcounter{section}{3}
\setcounter{subsection}{0}

\section{Extensões do Conceito de Limite}
\subsection{Limites no Infinito}
\subsection{Limites Infinitos}
\subsection{Sequência e Limite de Sequência}
\subsection{Limite de Função e Sequências}
\subsection{O Número $e$}

\end{document}
